% Results, wrapped around the five figures (rebuilt 2026-09-24 against the figure builds of that morning; the
% sixth figure, the LevelCheck comparison, moved to paper 2 that evening). First paragraph = the study population.
% Then the Methods items in order, under the same subheadings (REVIEW2_RSNA.md Fix 25): the phenotypes (Figure 1,
% Table 2), the level encoder (Figure 2), naming, counting and coupling on the intraoperative C-arm (Figure 4,
% Table 3 = tab3_naming.tex; detection in Table S2) and transitional sample size (Figure 5). Every percentage
% carries its denominator here or in the cited table; sources in NUMBERS.md ("Six-figure wrap" and "Revision 2").
% Where the numbers come from (REVISION_LOG2.md, all eval_gate HELD OUT): the zero-shot run is the lineage-complete
% v9 encoder + head (results/eval_gate/leaks/zeroshot_v9.json, overlap 0, all 148 cases); the fold-model, field,
% coupling and stop-signal numbers are the v9-lineage re-run on the PRE truth (results/eval_gate/closure/cv_v9/:
% cv_v9.json, cv_v9_extras_pre.json, coupling_pre/coupling.json; the file Figure 4 draws); the count from the sacrum
% and the detection counts use no level encoder (FIGREC:figure4); the embedding numbers are the v9 selection split
% and say so; the same-device learning curve is results/indomain/INVENTORY.md (HELD OUT rows, body boxes); 673
% agreed AP vertebrae in 143 cases (docs/carmlspine_cv5_mixed_2026-09-21.json#/oracle/AP/kfull), 631 in 142 detected.
% Revision 5 (2026-09-24, REVISION_LOG3.md section 8; PAPER1_SCOPE.md): the paragraph "The LevelCheck comparison"
% and the exact-test sentence of "Transitional sample size" moved verbatim to paper2_levelcheck/MOVED_FROM_PAPER1.md
% item 1; the sample-size paragraph quotes the pre-stated Wilson targets (35 persons at 90 %, 22 carriers at 85 %:
% Table 2 note f, results/power/power.json#n_needed_zero_errors, results/stats_review/recompute.json#R4_wilson). No
% transitional test person of the built split is scored here beyond the three C-ArmLSpine L6 cross-validation cases.
% Fill pass (2026-09-24 evening, REVISION_LOG3.md section 8 "Fill pass"): the study-population sentence quotes
% results/age_sex_by_source.json through \RESULT{agesex_*} (per-source detail in supplement Table S3). The earlier
% draft said the radiotherapy set releases no age or sex; COBRA2026 does (metadata.yaml, ct block: age on 115 and
% sex on 75 of the 145 detector cases), so the sentence now names only the intraoperative, cadaveric and phantom
% sets as releasing none.
% Data-case pass (2026-09-25, REVISION_LOG3.md section 9): the last placeholder (G5) is now the paragraph "Persons,
% films and coverage": film N curve (grid chain 40437811 -> results/coverage/film_ncurve/film_ncurve.json), the
% matched-N coverage arms (chain 40437814 -> results/coverage/experiment/scores.json) and the pooled nested fit
% (tools/indomain_fit_coverage.py -> results/coverage/pooled_fit.json, sensitivity/pooled_fit_conf074.json); every
% number is a \RESULT key of numbers_manifest.json (film_*, cov_*). The sample-size sentences now cite Figure 5a only
% (the old panels c, d moved to paper 2). The case-grouped lateral naming and the sacrum / body-box detection
% numbers of this file are \RESULT keys too (naming_lateral_cv_v9_*, det_carm_*); the rest stay typed until the
% RESULT-conversion pass.
% Acquisition-space pass (2026-09-25 afternoon, REVISION_LOG3.md section 10; Greg: "demonstrate how much they vary in
% the extra dimensions of the fluoro"): Figure 5 is now (a) persons needed, (b) naming vs training persons on film AND
% on the C-arm per view protocol (results/indomain/name_fill.json#figure5b), (c) same-device detection vs real frames
% with the one-view arms (results/indomain/det_fill.json curves.carm + protocol), (d) the matched-N render arms read
% along the lateral angle (experiment/runs/cov_surgio_b51_dr_*/unseen/recall_vs_angle.json). The paragraph "Persons,
% views and angles" walks the panels; the same-device sentences that closed "Naming, counting and coupling" moved into
% it; the pooled fit keeps its numbers without a panel. Keys name_carm_*, det_both_*, det_oneview_*, angle_*.
\section*{Results}

\paragraph{Study population.}
The corpus held 4,093 renders of labelled CT, 9,108 level-annotated films and 4,276 real exposures with
three-dimensional truth, 778 from mobile C-arms (Table~\ref{tab:data}). Of the intraoperative set's 910 cone-beam
CT vertebrae, 814 (89.5\,\%) were agreed, giving 676 agreed lateral vertebrae in 143 cases and 673
anteroposterior in 143. No age, sex or indication is released for the intraoperative, cadaveric or phantom sets.
Across the eight releases that give them (Table~S3), median age ran from \RESULT{agesex_csxa_age} years (CSXA) to
\RESULT{agesex_cobra_age} (COBRA2026) and the share female from \RESULT{agesex_ts_pct_female}\,\%
(TotalSegmentator) to \RESULT{agesex_csxa_pct_female}\,\% (CSXA).

\paragraph{Sources and transitional phenotypes.}
The sources held 194 transitional persons (Table~\ref{tab:transitional}): 110 with L6 (three intraoperative), seven
with T13, 28 with a sacralised L5 (12 with four lumbar vertebrae), 31 with Castellvi I/II, 14 with a lumbar rib and
four unassigned. No released radiograph or C-arm frame carries a human-read T13 or Castellvi grade
(Figure~\ref{fig:phenotypes}).

\paragraph{The level encoder and what it learns.}
On the checkpoint-selection split (13,009 crops of 1,612 persons), the nearest crop from another person was the
same level for 9,740 (74.9\,\%; 95\,\% confidence interval [CI]: 72.2, 77.6) (Figure~\ref{fig:embedding}a, b). L6
separated from L5 (silhouette 0.35; 0.31, 0.39; 98 crops, 10 persons; leave-person-out 0.34) more
than L2--L5 from the level above (0.16--0.23; L5 from L4, 0.23; 0.21, 0.25) and as far as L1 from T12 (0.32; 0.29,
0.35); the classifier named 75 of 98 L6 crops and 111 of 1,012 L5 crops L6 (Figure~\ref{fig:embedding}c; per
render source, supplement). In three
never-seen patients the contested L6 sat nearer the L6 cluster under whole-vertebra boxes only, and in two of four
fused-L5 validation persons the last free vertebra was named L5, the name a count gives (supplement). No T13 crop was outside training. Per unit area, occlusion
attribution was densest on the target body at each of six levels (enrichment 1.39--4.51), but for L4--L6 more fell on
the body below (36--43\,\%) than on the target (24--32\,\%) (Figure~\ref{fig:embedding}d).

\paragraph{Naming, counting and coupling on the intraoperative C-arm.}
The keypoint detector boxed 617 of 676 agreed lateral bodies (91.3\,\%), the held-out detector a sacrum box on 233
of the 286 frames showing one (81.5\,\%) (Figure~\ref{fig:exposure}a). Never trained on C-ArmLSpine, \netname{}
named 595 of 617 detected lateral vertebrae exactly (96.4\,\%; 92.3, 98.4) and 276 of 631 detected
anteroposterior (43.7\,\%), a paired difference of 54.3 points (45.6, 62.9; $P<.001$) on 593 vertebrae of 140 cases
(Table~\ref{tab:naming}). Fine-tuned by case-grouped folds, it named \RESULT{naming_lateral_cv_v9_kn} lateral (\RESULT{naming_lateral_cv_v9_pct}\,\%; \RESULT{naming_lateral_cv_v9_ci}) and
608 of 631 anteroposterior vertebrae (96.4\,\%; 92.5, 98.3), 1.0 point ($-1.9$, 4.0) and 52.6 points (44.1, 61.3)
over zero-shot; by acquisition-day folds, 568 of 617 lateral (92.1\,\%; 87.5, 96.0), 4.4 points below zero-shot
($-9.1$, 0.2). With three lateral bodies in view the case-grouped models named 1,079 of 1,159 (93.1\,\%; 89.2,
96.2), with one 451 of 603 (74.8\,\%) (Figure~\ref{fig:exposure}b); zero-shot, 970 of 1,159 (83.7\,\%) and 385 of 603
(63.8\,\%). On the 1,009 bodies of anchored frames (confidence 0.74), the count
from the sacrum was exact on 979 (97.0\,\%; 94.2, 99.2) (Figure~\ref{fig:exposure}c), the fine-tuned classifier on
986 (97.7\,\%) and the coupled call on 991 (98.2\,\%; 96.5, 99.5), 1.2 points over the count (0.0, 3.2) and 0.5
over the classifier ($-1.5$, 2.4). The read stopped on 23 of 233 anchored frames, catching five of
six erring frames and all four L6 calls: three on five-lumbar frames, one on the four anchored frames of the three
patients one labeller reads six-lumbar (Figure~\ref{fig:exposure}d).
The held-out detector found 231 of 287 vertebrae (80.5\,\%; 74.1, 86.8) and \RESULT{det_carm_sacrum_074_kn} sacra (\RESULT{det_carm_sacrum_074_pct}\,\%), short of the
pre-stated 97\,\% (Table~S2); on the \RESULT{denom_repaired_n} truth boxes that are complete, body-sourced and in the
field (supplement), \RESULT{det_heldout_repaired_kn} (\RESULT{det_heldout_repaired_pct}\,\%; \RESULT{det_heldout_repaired_ci}).
The radiotherapy imager's arc frames carried that detection: alone they found \RESULT{a3_e_only_carm_pct}\,\% of
C-ArmLSpine vertebrae, the cadaver C-arm's frames \RESULT{a3_d_only_carm_pct}\,\%, and adding them raised recall
\RESULT{a3_add_e_to_d_pp} points (\RESULT{a3_add_e_to_d_ci}) (Figure~S4).

\paragraph{Transitional sample size.}
With no error, 35 at-risk persons put the Wilson lower bound at the pre-stated 90\,\% and 22 carriers the
variant-sensitivity bound at 85\,\% (Figure~\ref{fig:power}a). No real-frame transitional test person was
verified: the three C-ArmLSpine L6 patients are scored only in cross-validation (Table~\ref{tab:naming}), and
against 91 unread COBRA2026 test-side screen calls (Table~\ref{tab:transitional}), a two-labeller count with both
junctions in the field found seven at-risk persons with an exported frame, whose error-free Wilson lower bound
would be 64.6\,\% (Figure~\ref{fig:power}a). The split's 51 at-risk CT test persons (most in the current
encoder's training data, Table~\ref{tab:transitional}) would give 93.0\,\% error-free.
\paragraph{Persons, views and angles.}
With a render-only encoder, exact-level naming on films rose with training persons (Figure~\ref{fig:power}b):
NHANES~II \RESULT{film_nhanes_acc_n5_pct} and \RESULT{film_nhanes_acc_all_pct}\,\% at five and all
\RESULT{film_nhanes_n_all_persons} persons; BUU-LSPINE \RESULT{film_buu_acc_n5_pct} and \RESULT{film_buu_acc_all_pct}\,\%
(\RESULT{film_buu_n_all_persons}); CSXA \RESULT{film_csxa_acc_n5_pct} and \RESULT{film_csxa_acc_all_pct}\,\%
(\RESULT{film_csxa_n_all_persons}). Transfer between corpora peaked at
\RESULT{film_transfer_nhanes_to_csxa_pct}\,\% (NHANES~II to CSXA, \RESULT{film_transfer_nhanes_to_csxa_n_films} films).
On the C-arm, naming separated by view rose like film: lateral from \RESULT{name_carm_lat_n0_pct}\,\% with no C-arm
person to \RESULT{name_carm_lat_n50_pct}\,\% with 50, anteroposterior from \RESULT{name_carm_ap_n0_pct}\,\% to
\RESULT{name_carm_ap_n5_pct}\,\% with five and \RESULT{name_carm_ap_nall_pct}\,\% with all; lateral-trained models named \RESULT{name_carm_lat_to_ap_nall_pct}\,\% of
anteroposterior bodies, and training both views cost neither (\RESULT{name_carm_both_lat_pct} and
\RESULT{name_carm_both_ap_pct}\,\%). Same-device detection (body boxes, 0.74) rose from
\RESULT{det_carm_body_recall_074_pct}\,\% (\RESULT{det_carm_body_recall_074_kn}) without the device's frames to a
plateau from 24 frames (lateral \RESULT{det_both_n24_lat_pct}, anteroposterior \RESULT{det_both_n24_ap_pct}\,\%;
\RESULT{det_both_n24_repaired_all_pct}\,\% on the repaired denominator); one-view arms of 25 or 50 frames found
\RESULT{det_oneview_own_min_pct}--\RESULT{det_oneview_own_max_pct}\,\% of its bodies and
\RESULT{det_oneview_other_min_pct}--\RESULT{det_oneview_other_max_pct}\,\% of the other's (Figure~\ref{fig:power}c).
On \RESULT{angle_n_subjects} unseen subjects' renders, render-only detectors at matched count
(\RESULT{cov_arm_surgio_matched_n} renders, \RESULT{cov_arm_surgio_subjects} subjects) trained on anteroposterior
and lateral poses alone found \RESULT{angle_low_trained_pct}\,\% at those angles and
\RESULT{angle_low_untrained_pct}\,\% over the \RESULT{angle_n_untrained_bins} other bins; doubling those frames gave
\RESULT{angle_low2n_trained_pct} and \RESULT{angle_low2n_untrained_pct}\,\%, the same count spread over the arc
\RESULT{angle_high_trained_pct} and \RESULT{angle_high_untrained_pct}\,\% (Figure~\ref{fig:power}d); the Elekta XVI
arms scored zero. On one patient's rotational acquisition the pooled detector found \RESULT{s5_pooled_spin_pct}\,\% of
bodies across the arc, the held-out detector \RESULT{s5_heldout_spin_pct}\,\% (Figure~S5). A nested fit
over \RESULT{cov_pooled_n_points} points preferred count plus lateral-angle coverage at confidence 0.25 (AICc lower by
\RESULT{cov_pooled_daicc}; $P\RESULT{cov_pooled_lr_p}$), not at 0.74 ($P\RESULT{cov_pooled_lr_p_conf074}$).
